\appendix

\section{Further background}
\label{app:further-background}

\subsection{Inner product spaces and spinors}
\label{sec:inner-products-spinors}

For the sake of brevity, we will not give a detailed account of the theory of Clifford algebras, spin groups and spinors here; standard references include \cite{Harvey1990,Lawson1989}; see also \cite{JMF-spin}. We will briefly establish some conventions, however. Throughout, let $(V,\eta)$ be a (pseudo-)inner product space.
We recall that using $\eta$, the vector spaces $V$ and $V^*$ can be identified via the \emph{musical isomorphisms}~${\flat\colon V\to V^*}$, $v\mapsto v^\flat$ where
$v^\flat(w) = \eta(v,w)
$
for all $v,w\in V$ and $\sharp\colon V^*\to V$, $\alpha\mapsto\alpha^\sharp$ which is defined as the inverse of $\flat$.

\subsubsection{Orthogonal groups and algebras}
\label{sec:orth-groups-algs}

We define the \emph{orthogonal group} of $(V,\eta)$ as the group of linear automorphisms of $V$ which preserve $\eta$,
\[
	\Orth(V,\eta) := \{A\in\GL(V) \mid \forall v,w\in V,\, \eta(Av,Aw) = \eta(v,w)\},
\]
and the \emph{special orthogonal group} as the subgroup of elements of $\Orth(V,\eta)$ which preserve orientations,
$\SO(V,\eta) := \{A\in \Orth(V) \mid \det A = 1\}$,
though we typically suppress the inner product $\eta$ in the notation where there is no ambiguity. The latter group is a normal subgroup of the former of index~2. Some basic topological properties depend on the signature of $\eta$ as follows.
\begin{itemize}\itemsep=0pt
	\item If $\eta$ is (positive- or negative-)definite, $\Orth(V)$ has two connected components, with the connected component of the identity being $\SO(V)$, and both groups are compact.
	\item If $\eta$ is not definite, $\Orth(V)$ has four connected components, two of which constitute $\SO(V)$. We denote the connected component of the identity by $\SO_0(V)$.\footnote{Following naming conventions used for the Lorentzian case in physics, this group is sometimes called the \emph{proper orthochronous orthogonal group}. It consists of invertible linear isometries of $(V,\eta)$ which preserve both orientations and time orientations on $V$.} None of these groups are compact.
\end{itemize}
Clearly, all of these groups share the same Lie algebra, which can be identified with the Lie algebra of $\eta$-skew-symmetric endomorphisms
\[
	\fso(V) := \{A\in\End(V) \mid \eta(Av,w) + \eta(v,Aw) = 0\}
\]
under the commutator bracket. We also define $\Orth(p,q):=\Orth(\RR^{p,q})$ etc., where $\RR^{p,q}$ is the standard inner product space with signature $(p,q)$; that is, with $p$ positive eigenvalues and $q$ negative eigenvalues.

\subsubsection{Clifford algebras}
\label{sec:clifford-algs}

The defining relation for the Clifford algebra $\Cl(V,\eta)$ is
\begin{equation}\label{eq:clifford-rel}
	v\cdot v = \pm\eta(v,v)\1
\end{equation}
for all $v\in V$. Note that there is a choice of sign convention here; the $+$ is more common in the physics literature while the $-$ is usually preferred in mathematics. We adopt the mathematical convention with the $-$ sign in defining $\Cl(p,q):=\Cl(\RR^{p,q})$. Any inner product space is (non-canonically) isomorphic to $\RR^{p,q}$ for some unique pair $(p,q)$, and such an isomorphism induces an isomorphism of associative algebras $\Cl(V,\eta)\cong \Cl(p,q)$. In previous works on supersymmetry and Spencer cohomology \cite{Beckett2021,deMedeiros2016,deMedeiros2018,Figueroa-OFarrill2017} and in the upcoming work treating 2-dimensional examples mentioned in the introduction, a $+$ sign is adopted for calculational convenience. We also typically suppress $\eta$ in the notation for the Clifford algebra (and the related groups defined below), writing $\Cl(V)=\Cl(V,\eta)$.

There is an isomorphism of vector spaces $\Wedge^\bullet V\simeq \Cl(V,\eta)$ which is most easily described in terms of an orthonormal basis $\{\be_\mu\}_{\mu=1}^{\dim V}$, although it is actually independent of this choice; to each exterior product of the form $\be_{\mu_1}\wedge\be_{\mu_2}\wedge\cdots\wedge\be_{\mu_r}$, we associate the Clifford product~${\be_{\mu_1}\cdot\be_{\mu_2}\cdot\cdots\cdot\be_{\mu_r}}$. We thus identify these two spaces and view Clifford multiplication $\cdot$ as a~new algebra structure on the exterior algebra $\Wedge^\bullet V$.

The exterior algebra $\Wedge^\bullet V$ has a natural $\ZZ$-grading by form degree. Clifford multiplication does not respect this grading, but it does respect the $\ZZ_2$-grading given by taking the $\ZZ$-grading modulo 2. We denote the even subalgebra of the Clifford algebra \big(which is generated by the subspace corresponding to $\Wedge^2 V$\big) by $\Cl_{\overline0}$.

There is a vector space embedding $\omega\colon\fso(V)\hookrightarrow\Wedge^\bullet V$ whose image is $\Wedge^2 V$ given by $A\mapsto \pm\frac{1}{2}\omega_A$ where $\omega_A$ is the 2-vector defined by the equation
$\imath_{v^\flat}\omega_A = -Av
$
where $\imath_{v^\flat}$ is contraction with the musical image $v^\flat\in V^*$ of an arbitrary vector $v\in V$.
This an embedding of Lie algebras if $\Wedge^\bullet V$ is equipped with the commutator Lie bracket with respect to Clifford multiplication.

\subsubsection{Pinor and spinor representations}
\label{sec:pinor-spinor}

The Clifford algebra $\Cl(V)$ has either a single isomorphism class of irreducible real finite-dimensional modules or two such classes, depending on $p-q\mod 4$ where $(p,q)$ is the signature of $(V,\eta)$.
\begin{itemize}\itemsep=0pt
	\item For $p-q\neq 3\mod 4$, there is a unique module which we denote by $\ssS$.
	\item For $p-q = 3 \mod 4$, there are two such modules which can be distinguished by the fact that a unit volume element in $\Wedge^{p+q}V$ acts as $+\1$ (the identity) on one of these modules and $-\1$ on the other; given an orientation on $V$, we denote by $\ssS$ the module on which the canonical unit volume element acts as $+\1$.
\end{itemize}
All of these are referred to as (irreducible, real) \emph{pinor} modules. Under the induced action of the even subalgebra $\Cl_{\overline 0}(V)$, or equivalently the spin group (see below), the pinor modules may be reducible, depending on $p-q\mod 8$.
\begin{itemize}\itemsep=0pt
	\item For $p-q=3,5,7 \mod 8$, $\ssS$ is irreducible; for later convenience we define $\ssS_1:=\ssS$ in this case.
	\item For $p-q=1,2,6$, $\ssS$ has a decomposition into irreducible $\Cl_{\overline0}(V)$-modules $\ssS=\ssS_+\oplus\ssS_-$ which happen to be isomorphic (as real representations); let us set $\ssS_1:=\ssS_+\cong\ssS_-$ here.
	\item For $p-q=0,4\mod 8$, there is a similar decomposition $\ssS=\ssS_+\oplus\ssS_-$, except here $\ssS_+\not\cong\ssS_-$. In any case, $\ssS_1$ or $\ssS_\pm$ are called the (irreducible, real) \emph{spinor} modules.
\end{itemize}

\subsubsection{Pin and spin groups}
\label{sec:pin-spin-group}

We define the \emph{pin group} and \emph{spin group} as the subgroups
\begin{gather*}
	\Pin(V) = \{v_1\cdot v_2\cdot\dots\cdot v_{r} \mid v_i\in V, \, \eta(v_i,v_i) = \pm 1 \},	\\
	\Spin(V) = \{v_1\cdot v_2\cdot\dots\cdot v_{2r} \mid v_i\in V, \, \eta(v_i,v_i) = \pm 1 \} = \Pin(V)\cap \Cl_{\overline0},
\end{gather*}
of the group of units $\Cl(V)^\times$. The spin group is a normal subgroup of the pin group of index 2, and similarly to the orthogonal groups, we have the following:
\begin{itemize}\itemsep=0pt
	\item If $\eta$ is definite, $\Pin(V)$ has two connected components, with the connected component of the identity being $\Spin(V)$, and both groups are compact.
	\item If $\eta$ is not definite, $\Pin(V)$ has four connected components, two of which constitute $\Spin(V)$. We denote the connected component of the identity by $\Spin_0(V)$. None of these groups are compact.
\end{itemize}
There is a canonical double-covering of Lie groups $\Pin(V)\to\Orth(V)$ which restricts to double-coverings $\Spin(V)\to \SO(V)$ and (where the groups exist) $\Spin_0(V)\to\SO_0(V)$. These covering maps differentiate to a canonical isomorphism of Lie algebras $\mathfrak{spin}_{\overline0}(V)=\mathfrak{spin}(V)\to\fso(V)$, thus we will identify these algebras.

The pin and spin groups act on the pinor modules by restriction of the action of the Clifford algebra. The irreducible pinor module $\ssS$ is irreducible under the action of $\Pin(V)$. It is irreducible under the action of $\Spin(V)$ or $\Spin_0(V)$ if and only if it is irreducible under $\Cl_{\overline 0}(V)$ and, in the appropriate signatures, $\ssS_1$ or $\ssS_\pm$ are irreducible under the action of these groups. The induced action of $\fso(V)$ on the (s)pinor modules is given by
\begin{equation}\label{eq:sov-spinor-action}
	A\cdot s = \pm\frac{1}{2}\omega_A\cdot s
\end{equation}
for $A\in\fso(V)$ and $s\in \ssS$, where the sign is the same as that in the Clifford relation \eqref{eq:clifford-rel} and the dot on the left-hand side is the Clifford action.

\subsection{Lie derivatives}
\label{sec:lie-der}

A standard concept in differential geometry, Lie derivatives play a central role in the definition of the Killing (super)algebra. In this appendix, we will first introduce some useful formulae for Lie derivatives of tensor fields and then show how these generalise to allow us to take Lie derivatives of spinor fields along Killing vectors.

\subsubsection{Lie derivative of vectors and tensors}
\label{sec:lei-der-vector-tensor}

Let $(M,g)$ be a connected pseudo-Riemannian manifold. We denote its Lie algebra of Killing vector fields\footnote{The notation should not be taken to imply that this is the Lie algebra of the isometry group of $(M,g)$ which is in general only a subalgebra of $\fiso(M,g)$; Killing vector fields may not integrate to a global action by a one-parameter group of isometries. We will not need to discuss the isometry group itself in this appendix or in the main text so this should not cause confusion.}
by $\fiso(M,g)$. The Levi-Civita connection will be denoted by $\nabla$. The bundle of skew-symmetric endomorphisms (with respect to $g$) of $TM$ will appear when we discuss Killing vectors. We denote the space of sections of this bundle by $\fso(M,g)$ and note that it is an (infinite-dimensional) Lie algebra under the commutator Lie bracket. The perspective presented here is due to Kostant \cite{Kostant1955}.

Given a vector field $X\in\fX(M)$, we define an endomorphism of $TM$ by
$A_X(Y) := -\nabla_Y X
$
for all $Y\in\fX(M)$. This gives us a useful formula for the Lie derivative along $X$: for $Y\in\fX(M)$,
\begin{equation}\label{eq:lie-der-vec-formula}
	\eL_X Y = \comm{X}{Y} = \nabla_X Y - \nabla_Y X = \nabla_X Y + A_X Y,
\end{equation}
where we have used the torsion-freeness of $\nabla$. We will also use
\begin{equation}\label{eq:comm-formula}
	\comm{X}{Y} = A_X Y - A_Y X.
\end{equation}
A similar formula holds a tensor field $T$ of any type
\begin{equation}\label{eq:lie-der-tensor-formula}
	\eL_X T = \nabla_X T + A_X \cdot T,
\end{equation}
where the action of $A_X$ on $T$ is defined via a Leibniz formula. Note that we have a different Leibniz rule of the form
\[
	\nabla_X (A_Y\cdot T) = \nabla_X(A_Y)\cdot T + A_Y\cdot\nabla_X T
\]
for all $X,Y\in\fX(M)$, or, abstracting $T$, $\comm{\nabla_X}{A_Y}=\nabla_X A_Y$ as endomorphisms of the space of tensor fields.

We can use the formula \eqref{eq:lie-der-tensor-formula} to show that $X$ is a Killing vector if and only if $A_X\in\fso(M,g)$: using the metric-compatibility of $\nabla$, we have $\eL_X g = A_X\cdot g$, so
\begin{equation}\label{eq:killing-condition}
	\eL_X g = 0 \iff g(A_X Y, Z) + g(Y,A_X Z) = 0,\qquad\forall Y,Z\in\fX(M).
\end{equation}
One can also show \cite[Lemma~2.2]{Kostant1955} that, for $X\in\fiso(M,g)$ and $Y\in\fX(M)$,
\begin{equation}\label{eq:nabla-A_X}
	\nabla_Y A_X = -R(Y,X),
\end{equation}
where $R$ is the Riemann curvature, so in particular $R(Y,X)\in\fso(M,g)$. Then for two Killing vectors $X,Y\in\fiso(M,g)$ and a vector field $Z\in\fX(M)$, using equations~\eqref{eq:comm-formula} and \eqref{eq:nabla-A_X} along with the algebraic Bianchi identity,
\begin{align*}
	A_{\comm{X}{Y}}Z
		&= -\nabla_Z\qty(A_X Y - A_Y X)= -(\nabla_ZA_X)Y - A_X\nabla_ZY + (\nabla_ZA_Y)X + A_Y\nabla_ZY\\
		&= R(Z,X)Y + A_XA_YZ - R(Z,Y)X- A_YA_XZ= \comm{A_X}{A_Y}Z - R(X,Y)Z.
\end{align*}
We thus have
\begin{equation}\label{eq:A-comm}
	A_{\comm{X}{Y}} = \comm{A_X}{A_Y} - R(X,Y)
\end{equation}
for all Killing vectors $X$ and $Y$.

\subsubsection{Spin structures and spinor bundles in arbitrary signature}
\label{sec:spin-structures}

Now let $(M,g)$ be a pseudo-Riemannian spin manifold of signature $(p,q)$. We recall that this means that $M$ is oriented and that there exists a $\Spin(p,q)$-principal bundle $P\to M$ over $M$ and a $\Spin(p,q)$-equivariant bundle map $\varpi\colon P\to F_{\rm SO}$, where $F_{\rm SO}\to M$ is the special orthonormal frame bundle of $(M,g)$ on which the spin group acts via the natural map $\Spin(p,q)\to\SO(p,q)$.

In indefinite signature, the structure group of the bundle $P\to M$ can be reduced to $\Spin_0(p,q)$ if and only if $(M,g)$ is time-orientable, and if this is the case then $\varpi$ restricts to a $\Spin_0(p,q)$-equivariant mapping $\varpi_0\colon P_0\to F_{SO_0}$ of the reduced bundles. Note that many authors assume time-orientability and refer to this reduced structure as the spin structure.\footnote{The reduced structure has been called a \emph{strong spin structure}, and manifolds admitting them \emph{strongly spin} \cite{Cortes2021,Shahbazi2024_2,Shahbazi2024_1}.}
	
Let $S$ be a (possibly $N$-extended) spinor module of $\Spin(p,q)$ in the sense discussed in Section~\ref{sec:isom-alg-spinors} and let $\Sbundle:=P\times_{\Spin(p,q)} S\to M$ be the associated vector bundle over $M$, which we will call the \emph{spinor bundle}. Sections of $\Sbundle$ will be called \emph{spinor fields}, and we denote the space of such fields by $\fS=\Gamma(\Sbundle)$. We also let $\Cl(M,g)=P\times_{\Spin(p,q)}\Cl(p,q)\cong F_{\rm SO}\times_{\SO(p,q)}\Cl(p,q)$ be the associated Clifford bundle (where the action is via the twisted adjoint, or the algebra automorphism induced by the natural action on $\RR^{p,q}$) and recall that we can identify it with the exterior bundle $\Wedge^\bullet T^*M$. Sections of this bundle can be Clifford-multiplied and Clifford-act on spinor fields via pointwise Clifford multiplication and action. The (spin-lift of) the Levi-Civita connection extends to all of these bundles, and there is a Leibniz rule for $\nabla$ with respect to Clifford multiplication and action.

The adjoint bundle $\ad F_{\rm SO}\cong\ad P$ can be identified with the space of $g$-skew-symmetric endomorphisms of $TM$, and it embeds in $\Cl(M,g)$ as $\Wedge^2 TM$. As a Lie algebra, the space of sections $\fso(M,g)$ acts on spinor fields via the Clifford action, i.e., we apply equation \eqref{eq:sov-spinor-action} pointwise.

\subsubsection{Lie derivative of spinors}
\label{sec:lie-der-spinors}

The formulae~\eqref{eq:lie-der-vec-formula} and \eqref{eq:lie-der-tensor-formula} suggest the following definition of the Lie derivative to spinor fields.

\begin{Definition}[spinorial Lie derivative \cite{Kosmann1971}]\label{def:lie-der}
	The Lie derivative of a spinor field $\epsilon\in\fS=\Gamma(\Sbundle)$ along a Killing vector $X\in\fiso(M,g)$ is given by
$
		\eL_X \epsilon := \nabla_X \epsilon + A_X\cdot \epsilon$.
%	\label{eq:lie-der-spin-formula}
\end{Definition}

Note, however, that while the formulae for the Lie derivatives of vector and tensor fields hold for the Lie derivative along any vector field $X$, the right-hand side of the spinor formula is only defined for Killing vectors\footnote{The spinorial Lie derivative is actually defined for \emph{conformal Killing} vector fields in \cite{Kosmann1971}, but we will only use it for Killing vectors here.}
$X$, since sections of $\End(TM)$ do not act on $\fS$, while~$\fso(M,g)$ does. Note that $\fso(M,g)$ is an (infinite-dimensional) Lie algebra and the action on $\fS$ defines a~representation. This action obeys a Leibniz rule of the form $\comm{\nabla_X}{A}=\nabla_X A$ as endomorphisms of $\fS$ for all $A\in\fso(M,g)$ and $X\in\fX(M)$; more explicitly,
$\nabla_X (A\cdot\epsilon) = (\nabla_X A)\cdot\epsilon + A\cdot(\nabla_X \epsilon)
$
for all $\epsilon\in\fS$. It follows immediately from the above that we have yet another Leibniz rule~${\eL_X (A\cdot\epsilon) = (\eL_X A)\cdot\epsilon + A\cdot(\eL_X \epsilon)
}$
for all $X\in\fiso(M,g)$, $A\in\fso(M,g)$, or equivalently~${\comm{\eL_X}{A}=\eL_X A}$ as endomorphisms of $\fS$. One can also show that where $X$ is a Killing vector field and $Y$ any vector field, the following compatibility condition between the spinorial Lie derivative and the Levi-Civita connection on spinors holds
\begin{equation}\label{eq:lie-der-compat}
	\comm{\eL_X}{\nabla_Y} = \nabla_{\comm{X}{Y}},
\end{equation}
while if both $X$ and $Y$ are Killing vector fields,
$\eL_{\comm{X}{Y}} = \eL_X\eL_Y - \eL_Y\eL_X$.
In particular, the spinorial Lie derivative gives $\fS$ the structure of an $\fiso(M,g)$ module.

We can also use the Leibniz rule to define an action of $\fso(M,g)$, and thus a Lie derivative along Killing vectors, on sections of $\End\Sbundle$ and on spaces of forms with values in this bundle~$\Omega^p(M;\End\Sbundle)$. More generally, the space of sections of any bundle constructed by taking duals and tensor products of $TM$ and $\Sbundle$ is equipped with such a structure, and all the Leibniz rules and compatibility conditions discussed above also hold when acting on these spaces.

%%%%%%%%%%%%%%%%%%%%%%%%%%%%%%%%%%%%%%%%%%%%%%%%%%%%
%%%%%%%%%%%%%%%%%%%%%%%%%%%%%%%%%%%%%%%%%%%%%%%%%%%%

\subsection{Grading and filtration of superspaces and superalgebras}
\label{sec:grading-filtration-superspace}

We assume that the reader is familiar with ordinary $\ZZ$-graded vector spaces and with vector superspaces but recall and set some terminology for $\ZZ$-graded superspaces and algebras, and we also introduce filtrations.

\subsubsection{Graded vector superspaces and superalgebras}
\label{sec:graded-superalg}

A vector superspace $V$ is said to be \emph{$\ZZ$-graded} with grading $V=\bigoplus_{k\in\ZZ}V_k$ if the latter is a~$\ZZ$-grading of $V$ as a vector space and the $\ZZ$-grading is compatible with the $\ZZ_2$ grading on $V$ in the sense that an element of degree $d$ in the $\ZZ$-grading has degree (parity) $\overline{d}=d\mod 2$ in the $\ZZ_2$ grading. Note that we can always make a $\ZZ$-graded \emph{vector space} into a $\ZZ$-graded vector \emph{superspace} by taking the $\ZZ$-grading modulo $2$. We say that a $\ZZ$-graded vector (super)space~${V=\bigoplus_{k\in\ZZ}V_k}$ has \emph{finite depth} $h\in\ZZ_{\geq 0}$ if $V_{-h}\neq 0$ but $V_k=0$ for all $k<-h$.

A Lie superalgebra $\fg$ is said to be \emph{$\ZZ$-graded} with grading $\fg=\bigoplus_{k\in\ZZ}\fg_k$ if the latter is a~$\ZZ$-grading of $\fg$ as a vector superspace and the superbracket respects this grading in the sense that~${\comm{\fg_k}{\fg_l}\subseteq\fg_{k+l}}$.

If $V$ is a vector superspace, we define the exterior algebra $\Wedge^\bullet V:=\mathcal{T}(V)/I$ where $\mathcal{T}(V)$ is the tensor algebra and I is the ideal generated by elements of the form $v\otimes w + (-1)^{\abs{v}\abs{w}}w\otimes v$, where~$v$,~$w$ are homogeneous elements with parities $\abs{v}$, $\abs{w}$, respectively. This is a $\ZZ$-graded associative algebra of depth 0: the grading is $\Wedge^\bullet V = \bigoplus_{k=0}^\infty \Wedge^k V$ where $\Wedge^k V$ is the image of~${\mathcal{T}^k(V)= V^{\Otimes k}}$ under the quotient map. One can show that, for $V=V_{\overline{0}}\oplus V_{\overline{1}}$,
\begin{equation}\label{eq:super-wedge}
	\Wedge^k V = \bigoplus_{i+j=k}\Wedge^i V_{\overline{0}}\otimes\Odot^j V_{\overline{1}}.
\end{equation}
It will occasionally be useful to remind the reader of this definition of $\Wedge^k$ of a vector superspace; when we do so, we will say that we mean it in the ``super-sense''. There is of course an analogous ``super'' version of the symmetric algebra $\Odot^\bullet V$ but we will not require it in this work.

\subsubsection{Filtered vector spaces and algebras}
\label{sec:filtered-alg}

We say that a vector space $V$ is \emph{$\ZZ$-filtered} if there exists a decreasing sequence of vector subspaces~${V \supseteq \dots \supseteq V^{k-1} \supseteq V^k \supseteq V^{k+1} \supseteq \dots \supseteq 0}$.
To emphasise that we are taking $V$ with its filtration, we will denote it by~$V^\bullet$. We say that~$V^\bullet$ is \emph{finite} if there exist some $k_1$, $k_2$ such that~${V^{k_1}=V}$ and $V^{k_2}=0$. Of course, if $V$ is finite-dimensional then $V^\bullet$ must be finite. We say that $V^\bullet$ (not necessarily finite) has \emph{finite depth}~${h\in\ZZ_{\geq 0}}$ if $V_{-h} = V$ but $V_k\neq V$ for all $k>-h$.

Given a $\ZZ$-filtered vector space, there is an \emph{associated graded vector space}
\[
	\Gr V^\bullet = \bigoplus_{k\in\ZZ} \Gr_k V,\qquad \text{where}\quad \Gr_k = \faktor{V^k}{V^{k+1}}.
\]
If the filtration $V^\bullet$ is finite then so is the grading on $\Gr V^\bullet$, and if $V^\bullet$ has finite depth $h$, so does $\Gr V^\bullet$. The associated graded vector space $\Gr V^\bullet$ is (non-canonically) isomorphic to~$V^\bullet$, which can be seen as follows. For each $V^k$, we choose a complement $W_k$ to $V^{k+1}$ so that~${V^k=W_k\oplus V^{k+1}}$. Then we have a grading $V=\bigoplus_{k\in\ZZ} W_k$, and the quotient map $V^k\to\Gr_k V^\bullet$ restricts to a linear isomorphism $W_k\to\Gr_kV^\bullet$. Taking the sum of these maps gives us a graded vector space isomorphism $V\to\Gr V^\bullet $.

If $V^\bullet$ is equipped with a Lie algebra structure, then it is \emph{$\ZZ$-filtered as an algebra} if $\comm{V^k}{V^l}\subseteq V^{k+l}$. Its associated graded vector space is naturally a $\ZZ$-graded Lie algebra; for $X\in \Gr_k V^\bullet, Y\in \Gr_l V^\bullet$, we choose representatives $\tilde{X}\in V^k,\tilde{Y}\in V^l$ and then define a bracket $\comm{-}{-}_{\mathrm{Gr}}$,
\[
	\comm{X}{Y}_{\mathrm{Gr}} := \overline{\big[\tilde{X},\tilde{Y}\big]},
\]
where the overline denotes the natural projection $V^{k+l}\to\Gr_{k+l}V^\bullet$. One can check that this bracket is well-defined precisely because of the condition that the algebra operation on $V^\bullet$ respects the filtration, it satisfies the appropriate axioms, and it is graded by construction.

We note that if $V=\bigoplus_{k\in\ZZ}V_k$ is a graded Lie algebra then it possesses a natural filtration~${V^k=\bigoplus_{l=k}^\infty V_l}$ which has finite depth $h$ if and only if the grading does. Note that we have~${V^k=V_k\oplus V^{k+1}}$, and there is a canonical isomorphism of graded vector spaces $V\to\Gr V^\bullet$ defined as for a general choice of complements $W_k$ as described above. In fact, it is easy to see that this is an isomorphism of graded algebras.

A \emph{filtered morphism} $f\colon V^\bullet\to W^\bullet$ of filtered vector spaces is a vector space morphism~$f\colon V\to W$ such that $f\bigl(V^i\bigr)\subseteq W^i$. We analogously define a \emph{filtered morphism} of filtered Lie superalgebras. A filtered morphism induces an associated graded morphism $\Gr f\colon \Gr V^\bullet \to W^\bullet$ by~${\Gr f([v])=[f(v)]}$. One can easily check that this is well-defined and respects gradings. We say that a filtered morphism is \emph{strict} if $f\bigl(V^k\bigr)=f(V)\cap W^k$.

If $f$ is an injective (respectively surjective) strict filtered morphism, then $\Gr f$ is injective (respectively surjective). Moreover, a linear isomorphism which is filtered is strict if and only if its inverse is filtered, and the inverse of a strict filtered isomorphism is also strict.\footnote{Analogous results do not hold for non-strict filtered morphisms, hence the need to introduce this technicality. See, e.g., \cite[\href{https://stacks.math.columbia.edu/tag/0108}{Example 0108}]{stacks-project}.}
We say that two filtered vector spaces (resp.\ Lie superalgebras) are \emph{isomorphic} if there exists a \emph{strict} filtered isomorphism between them.

\subsubsection{Filtered vector superspaces and superalgebras}
\label{sec:filtered-superspace-superalg}

If $V=V_{\overline{0}}\oplus V_{\overline{1}}$ is a vector superspace, we say that a vector space filtration $V \supseteq \dots \supseteq V^{k-1} \supseteq V^k \supseteq V^{k+1} \supseteq \dots \supseteq 0$ (or $V^\bullet$ for short) is a filtration of the superspace if it is compatible with the $\ZZ_2$-grading in the following sense. We demand that the filtered components are graded subalgebras
\begin{equation}\label{eq:filter-super-1}
	V^k = V^k_{\overline{0}} \oplus V^k_{\overline{1}},
\end{equation}
(note that $V^\bullet_{\overline{0}}$ and $V^\bullet_{\overline{1}}$ are filtrations of the even and odd subspace respectively) such that
\begin{equation}\label{eq:filter-super-2}
	V^{2k-1}_{\overline{0}} = V^{2k}_{\overline{0}}
	\qquad	\text{ and } \qquad
	V^{2k}_{\overline{1}} = V^{2k+1}_{\overline{1}}.
\end{equation}
Note that this means $\Gr_{2k+1}V^\bullet_{\overline 0}=\Gr_{2k}V^\bullet_{\overline 1}=0$, while $\Gr_{2k} V^\bullet=\Gr_{2k} V^\bullet_{\overline0}$ and $\Gr_{2k+1} V^\bullet = \Gr_{2k+1} V^\bullet_{\overline 1}$. This ensures that one can define a natural $\ZZ_2$-grading on $\Gr V^\bullet$ compatible with the $\ZZ$-grading.\footnote{These compatibility conditions seem to be somewhat under-emphasised in the literature; the need for such conditions is mentioned in the classic reference \cite{Kac1977}, but sufficient conditions are only given for a particular class of filtrations there, not for a general filtration. Our main reference \cite{Cheng1998} also mentions that a compatibility condition is required but does not spell it out. The conditions given here appear in \cite{Scheunert1979}.}

If $V^\bullet$ is a Lie superalgebra which is $\ZZ$-filtered as a vector superspace, then it is \emph{$\ZZ$-filtered as a Lie superalgebra} if $\comm{V^k}{V^l}\subseteq V^{k+l}$. One then defines a Lie superalgebra bracket on the associated graded superspace $\Gr V^\bullet$ in the obvious way.

To avoid overloading notation, in the main text we often drop the $\bullet$ superscript from filtered vector spaces and the $\mathrm{Gr}$ subscript from the operation on graded algebras; the filtration structure should be implicitly understood whenever we discuss such spaces. We note that a bullet superscript is also used for cohomological grading elsewhere in this work.
	
\subsection{Homogeneous spaces}
\label{sec:homog-spaces}

Here we establish some terminology, conventions and important results for working with homogeneous spaces. Such spaces play an important role in this work due to the homogeneity theorem (see Theorem~\ref{thm:homogeneity}): a highly supersymmetric supergravity background is (at least locally) homogeneous.

We give our own presentation of standard material here, but some authoritative references for material with no explicit citation in the text are the books of Kobayashi--Nomizu \cite{Nomizu1969} and Sharpe \cite{Sharpe1997}. Wise's work \cite{Wise2010} on the Cartan-geometric formulation of gravity also contains a~brief pedagogical introduction to (metric) Klein pairs.
	
\subsubsection{Homogeneous spaces and Klein pairs}

\begin{Definition}
	Let $G$ be a Lie group. A \emph{homogeneous $G$-space} is a manifold $M$ on which $G$ acts smoothly and transitively. A morphism of such spaces is a $G$-equivariant smooth map, and an isomorphism is a $G$-equivariant diffeomorphism.
	The \emph{isotropy group} of a point $p$ in such a~space is the stabiliser subgroup $G_p=\{g\in G \mid {g\cdot p = p} \}$.
\end{Definition}

Let $G$ be a Lie group and $M$ a homogeneous $G$-space. Fixing a point $p\in M$, we note that the \emph{orbit map} $\alpha_p\colon G\to M$ given by $\alpha_p(g) = g\cdot p$ is surjective (by transitivity), and $G_p=\alpha_p^{-1}(p)$, so~$G_p$ is a closed subgroup of $G$. If $G$ is connected, $M$ must also be connected since it is the image of $G$ under $\alpha_p$. We will often take $G$ to be connected when discussing homogeneous $G$-spaces. Note that all of the isotropy groups of the action of $G$ on $M$ are conjugate; for all~${g\in G}$ we have $G_{g\cdot p} = gG_p g^{-1}$.

%%%%%%%%%%% Isotropy representation and Frobenius reciprocity

The derivative at $p\in M$ of the diffeomorphism $\phi_g\colon M\to M$ associated to an element $g\in G$ is an invertible linear map $d_p\phi_g\colon T_pM\to T_{g\cdot p}M$; in particular, if $h\in G_p$ then $d_p\phi_h\in \GL(T_pM)$. This defines a representation of $G_p$ on $T_pM$ known as the \emph{isotropy representation}.

The derivative of $\alpha_p$ at the identity of $G$ is a linear map $d_e\alpha_p\colon \fg=T_eG\to T_pM$ with kernel~${\fg_p=\Lie G_p}$. One can show that $d_e\alpha_p\circ\Ad_h=d_p\phi_h\circ d_e\alpha_p$ for all $h\in G_p$; in other words, $d_e\alpha_p$ is $G_p$-equivariant where $G_p$ acts on $\fg$ via the restriction of the adjoint representation of $G$. Thus there exists a short exact sequence of $G_p$-modules
\[
\begin{tikzcd}
	&0 \ar[r] &\fg_p \ar[r] & \fg \ar[r,"d_e\alpha_p"] & T_pM \ar[r] & 0.
\end{tikzcd}
\]

The isotropy representation is of fundamental importance in the geometry of homogeneous spaces due to a particular manifestation of \emph{Frobenius reciprocity}: $G$-invariant tensor fields on~$M$ are in one-to-one correspondence with $G_p$-invariant tensors on $T_p M$. This correspondence is easy to describe: if $\tau$ is a $G$-invariant tensor field then its value $\tau_p$ at $p$ is clearly a $G_p$-invariant tensor on $T_p M$; conversely, if $t$ is a $G_p$-invariant tensor on $T_p M$ then we can define a tensor field $\tau$ on $M$ by defining $\tau_{g\cdot p}:=(d_p\phi_{g})(t)$. This tensor field $\tau$ is well-defined and $G$-invariant precisely because $t$ is $G_p$-invariant and $d\phi_{gg'}=d\phi_g\circ d\phi_{g'}$. We will see a more general version of Frobenius reciprocity in due course (see Proposition~\ref{thm:Frobenius-recipr}).

\begin{Definition}\label{def:lie-pair-algs}
	A \emph{Lie pair} is a pair $(\fg,\fh)$ where $\fg$ is a~finite-dimensional Lie algebra and $\fh$ a~subalgebra of $\fg$.
	A \emph{Klein pair} is a pair $(G,H)$ where $G$ is a connected Lie group and $H$ is a~closed subgroup of~$G$.
	The \emph{homogeneous $G$-space associated to the Klein pair $(G,H)$} is the left coset space $G/H$.
\end{Definition}

If $(G,H)$ is a Lie pair, the coset space $G/H$ is indeed a homogeneous $G$-space; clearly $G$ acts on it transitively, and it is a smooth manifold because $H$ is closed. It is connected because $G$ is connected, and $H$ is the isotropy group for $o:=H$ (the coset of the identity). Note that we assume that $G$ is connected but $H$ may not be connected. We will discuss the effect of this on the topology of the associated homogeneous space in Appendix~\ref{sec:homog-spaces-principal-bundles}.

Conversely, if $M$ is a homogeneous space for a connected group $G$ and $p\in M$ then $(G,G_p)$ is a Klein pair, and $M$ is of course the associated homogeneous space up to isomorphism of $G$-spaces; $M\cong G/G_p$. Fixing a connected Lie group $G$, one can show that there is a one-to-one correspondence between conjugacy classes of closed subgroups of $G$ and isomorphism classes of homogeneous $G$-spaces.

The discussion of the isotropy representation above shows that given a Lie pair $(G,H)$, we have $T_o(G/H)\cong \fg/\fh$ as $H$-modules, and that $H$-invariant tensors on $\fg/\fh$ induce $G$-invariant tensor fields on $G/H$. We will be particularly interested in $H$-invariant (pseudo-)inner products~$\eta$ on $\fg/\fh$, which induce $G$-invariant pseudo-Riemannian metrics $g$ on $G/H$.

\begin{Definition}A \emph{metric Lie pair} is a triple $(\fg,\fh,\eta)$ where $(\fg,\fh)$ is a Lie pair and $\eta$ is an $\fh$-invariant inner product on $\fg/\fh$.
A \emph{metric Klein pair} is a triple $(G,H,\eta)$ where $(G,H)$ is a~Klein pair and $\eta$ is an $H$-invariant inner product on $\fg/\fh$.
The \emph{homogeneous pseudo-Riemannian $G$-space associated to $(G,H,\eta)$} is the pair $(G/H,g)$, where $g$ is the pseudo-Riemannian metric induced by $\eta$.
\end{Definition}

Given a metric Klein pair $(G,H,\eta)$, set $M=G/H$, $o=H$ and note that since $\eta$ is $H$-invariant, the image of the isotropy representation, which we now denote by $\varphi\colon H\to GL(T_oM)$, is contained in $\Orth(T_oM)\subseteq\GL(T_oM)$. If $M$ is oriented, then since $G$ is connected, its action is orientation-preserving, hence the action of $H$ on $T_oM$ must also be orientation-preserving, so~${\Im\varphi\subseteq\SO(T_oM)}$. If the signature of $\eta$ is indefinite and $M$ is both oriented and time-oriented, we similarly find that $\Im\varphi\subseteq\SO_0(T_oM)$.

\subsubsection{Homogeneous spaces as principal bundles and their cohomology}
\label{sec:homog-spaces-principal-bundles}

%%%%%%%%%%% Principal bundles and connectivity

Given a Klein pair $(G,H)$, right multiplication by elements of $H$ on $G$ gives the coset map $G\to M=G/H$ ($g\mapsto gH$) the structure of a (right) $H$-principal bundle over $M$. Note that we \emph{only} consider the \emph{left} action of $G$ on itself, and $G$ does not act on the right of $M$ in a natural way, while $H$ acts on both $G$ and $M$ from the left and right -- the left action on $M$ is non-trivial, while the right action is trivial. The coset map is equivariant with respect to both the left action of $G$ and the right action of $H$.

We have a fibration $H\to G\to M$, hence a long exact sequence of groups in homotopy ending~in
\[
\begin{tikzcd}
	\cdots \ar[r] &\pi_1(G) \ar[r] &\pi_1(M) \ar[r] &\pi_0(H) \ar[r] &\pi_0(G)=1,
\end{tikzcd}
\]
where we take the basepoints of $G$ and $H$ to be the identity element and that of $M$ to be $o=H$. The group structure on $\pi_0(H)$ is induced by that of $H$, and $\pi_0(G)$ is trivial since we assume $G$ to be connected. Thus if $M$ is simply connected, $H$ must be connected. The converse is not true in general, but if we assume that $G$ is simply connected, $\pi_1(M)\cong\pi_0(H)$, so $M$ is simply connected if and only if $H$ is connected.

When working with a homogeneous $G$-space $M=G/H$, it is often useful to pass to the universal cover $\widetilde{G}$ of $G$ by pulling back the action of $G$ on $M$ to $\widetilde{G}$. We can then consider $M$ as a homogeneous space corresponding to the Klein pair \smash{$\bigl(\widetilde{G},H'\bigr)$} where $H'$ is the preimage of $H$ under the covering map \smash{$\widetilde{G}\to G$}. Note that the universal cover $\widetilde{M}$ of $M$ is then a homogeneous space corresponding to $\bigl(\widetilde{G},H'_0\bigr)$, where $H'_0$ is the connected component of $H'$.

\subsubsection{Associated bundles and Frobenius reciprocity}

For the remainder of this section, we let $V=T_o M$ for our convenience. Many natural bundles over~$M$ can be described as associated bundles with respect to representations of $H$. For example, there is a $G$-equivariant isomorphism of vector bundles
$
	TM \cong G\times_H V$,
where $H$ acts on $V$ via the isotropy representation and $G$ acts on $TM$ via the push-forward of the action on $M$ and on the associated bundle via $g\cdot[g',v]=[gg',v]$. The isomorphism in the direction~${G\times_H V\to TM}$ is given by $[g,v]\mapsto g_* v$. The isotropy representation also induces representations of $H$ on~${V^*=T^*_o M}$ and its exterior powers, giving us further $G$-equivariant isomorphisms
$T^*M \cong G\times_H V^*$,	$
	\Wedge^\bullet T^*M \cong G\times_H \Wedge^\bullet V^*$.
The action of $H$ on $\End(V)$ by conjugation by elements of $\GL(V)$ similarly gives us
$\End(TM) \cong G\times_H \End(V)$.

Considering the isotropy representation as a Lie group morphism $H\to\GL(V)$ gives various actions of $H$ on $\GL(V)$ -- by conjugation or by left or right multiplication by the images of elements in $H$, for example -- but we will concentrate on the left action by left multiplication. This allows us form a (right) $\GL(V)$-principal bundle $G\times_H \GL(V)$ equipped with a compatible (left) action of $G$, where the two groups act by right multiplication on the right factor and left multiplication on the left factor respectively. Now, fix a frame $f=(f_1,\dots,f_n)$ at $o$ ($n=\dim M$), that is, $f$ is an ordered basis for $V$. This choice induces a group isomorphism $\GL(V)\cong\GL(\RR^n)$ by representing elements of $\GL(V)$ in the frame $f$; in Einstein notation, for $A\in\GL(V)$ we define the components of the matrix $\underline{A}\in \GL(\RR^n)$ by $A(f_i)=A\indices{^j_i}f_j$. This allows us to treat the frame bundle $FM\to M$ as a $\GL(V)$-principal bundle; if $f'\in FM$ is another frame (at any point), we have
\smash{$(f'\cdot A)_i = f'_jA\indices{^j_i}$}.
This choice of frame $f$ also allows us to define a $G$-equivariant isomorphism of $\GL(V)$-principal bundles (that is, a $(G,\GL(V))$-equivariant bundle morphism)
$FM \cong G\times_H \GL(V)$,
where the left action of $G$ on frames is induced by its action on tangent vectors,
\[%\label{eq:g-f}
	g\cdot (f'_1,\dots,f'_n) = (g\cdot f'_1,\dots,g\cdot f'_n) = (g_*f'_1,\dots,g_*f'_n);
\]
the isomorphism is given by mapping $G\times_H \GL(V)\ni [g,A]\mapsto g\cdot f \cdot A$.

Given a \emph{metric} Klein pair $(G,H,\eta)$ with orientation, the isotropy representation $H\to\SO(V)$ allows us to construct a $G$-equivariant isomorphism of $\SO(V)$-bundles
$F_{\rm SO} \cong G\times_H \SO(V)$,
where~${F_{\rm SO}\to M}$ is the special orthonormal frame bundle of $(M,g)$, in a completely analogous way -- the construction requires a choice of \emph{oriented orthonormal} frame $f$. Similarly, without an orientation we have $F_O\cong G\times_H\Orth(V)$, and with time orientation we have $F_{SO_0}\cong G\times_H\SO_0(V)$, where the notation should be self-explanatory.

We can make particularly fruitful use of these isomorphisms with Frobenius reciprocity, which we have already seen an application of. This a powerful result which appears in a number of different guises in the literature but in general relates representations of a group $G$ restricted to a subgroup $H$ to representations of $G$ induced by representations of $H$. The most useful version to us here is due to Bott.

\begin{Proposition}[Frobenius reciprocity theorem \cite{Bott1965}]\label{thm:Frobenius-recipr}
	Let $(G,H)$ be a Klein pair, let $W$ be a~$G$-module and let $W'$ be an $H$-module. Then there is a an isomorphism of vector spaces
	\[
		\Hom_G(W,\Gamma(G\times_H W'\to G/H)) \simeq \Hom_H(W,W').
	\]
\end{Proposition}

In the statement above, we treat $W$ as an $H$-module by restricting the action of $G$ on $W$ and we treat the space of sections $\Gamma(G\times_H W'\to G/H)$ as a $G$-module by acting on the left factor of each fibre. We recover our previous statement about invariant tensor fields by noting, for instance, that, setting $W=\RR$ the trivial module and $W'=V=T_oM$, the $G$-equivariant bundle isomorphism $TM \cong G\times_H V$ gives
\smash{$\fX(M)^G \simeq V^H$},
so $G$-invariant vector fields on $M$ are in correspondence with $H$-invariant vectors at $o$. Similarly, using the $G$-equivariant bundle morphism $G\times_H \Wedge^\bullet V^* \cong \Wedge^\bullet T^*M$,
we have
\smash{$\Omega^\bullet(M)^G \simeq (\Wedge^\bullet V^*)^H$},
and similarly $\End(TM)^G\simeq \End(V)^H$ etc.

\subsubsection{Connections, Wang's theorem and Nomizu maps}
\label{sec:nomizu}

Just like tensor fields, invariant connections on homogeneous spaces have a simpler description in terms of equivariant linear maps. This formalism was introduced for affine connections on the homogeneous spaces of reductive Lie pairs by Nomizu \cite{Nomizu1954}. We will discuss a generalisation due to Wang \cite{Wang1958}. A good reference for all statements is \cite{Nomizu1969}.

For a Lie pair $(G,H)$, we once again denote $M=G/H$, $V=T_oM\cong\fg/\fh$ where $o=H$. Note that since $H$ acts on $V$ via the isotropy representation which we denote $\varphi\colon H\to\GL(V)$, it acts on $\fgl(V)$ by conjugation.

%%%%%%%%%%% Wang's theorem

Suppose that $\pi\colon P\to M$ is a right principal $K$-bundle for some Lie group $K$, and suppose that the (left) action of $G$ on $M$ lifts to a (left) action on $P$ which is compatible with the action of $K$; that is, we have an action of $G$ on $P$ such that $\pi$ is $G$-equivariant and $(g\cdot p)\cdot k=g\cdot (p\cdot k)$ for all $g\in G$ and $k\in K$. Then note that $H$ preserves the fibre $P_o$, and by fixing a~basepoint~${p\in P_0}$ we can define a~group homomorphism $\psi\colon H\to K$ by declaring $h\cdot p = p\cdot\psi(h)$. A~different choice of basepoint changes $\psi$ by conjugation by an element of $k$. We can now state the following.\looseness=1

\begin{Theorem}[\cite{Wang1958}]\label{thm:Wang}
	Let $(G,H)$ be a Lie pair and let $P\to M=G/H$ be a right principal $K$-bundle for some Lie group $K$. Suppose that the left action of $G$ lifts to $P$ compatibly with the action of $K$. Fix a basepoint $p\in P_o$ and let $\psi\colon H\to K$ be the induced group morphism. Then there is a one-to-one correspondence between~$G$-invariant principal connections $\omega\in\Omega^1(P;\fk)$ on~$P$ and linear maps $\Psi\colon \fg\to\fk$ such that
	\begin{enumerate}\itemsep=0pt
		\item[$(1)$] $\Psi\circ\Ad^G_h = \Ad^K_{\psi(h)}\circ\Psi$ for all $h\in H$,
		\item[$(2)$] $\Psi(X)=(d_e\psi)(X)$ for all $X\in\fh$.
	\end{enumerate}
	The correspondence is given by
$\Psi(X) = \omega_p(\xi_X)
$
	where $X\in \fg$ and $\xi_X$ is the fundamental vector field on $P$ associated to $X$. Moreover, the curvature $F_\omega\in\Omega^2(P;\fk)^G$ is given by
	\[
(F_\omega)_p(\xi_X,\xi_Y) = \comm{\Psi(X)}{\Psi(Y)}_\fk - \Psi(\comm{X}{Y}_\fg).
	\]
\end{Theorem}

The final equation uniquely determines $\omega_p$ since $G$-transitivity on the base implies that any fibre is carried to $P_o$ by the action of some element of $G$, $K$ acts freely and transitively on each fibre, and $\omega$ is a $G$-invariant principal connection.

We will refer to the map $\Psi$ as the \emph{Nomizu map}, since it plays the same role in Wang's theorem that the more well-known notion of Nomizu map plays in the reductive case due to Nomizu.

If $P=FM$ is the frame bundle, recall from Appendix~\ref{sec:homog-spaces-principal-bundles} that a choice of frame~$f$ in~$F_oM$ gives~$FM$ the structure of a principal $K=\GL(V)$-bundle and that we have a compatible left action of~$G$. Then we immediately have $\psi=\varphi\colon H\to\GL(V)$, the isotropy representation. Moreover, if $(G,H,\eta)$ is an orientable metric Lie pair, we can similarly take $P=F_{\rm SO}$, ${K=\SO(V)}$ and $\psi=\varphi\colon H\to\SO(V)$. Now recall that an affine connection is equivalent to a principal connection on $FM$, so Wang's theorem tells us that such a connection on a homogeneous space is characterised by a map $\Psi\colon \fg\to\fgl(V)$, and similarly, metric affine connections on a homogeneous pseudo-Riemannian space correspond to maps $\Psi\colon \fg\to\fso(V)$ (subject to the conditions of the theorem).

We finish this section by describing the Nomizu map associated to the Levi-Civita connection for a metric Lie pair $(G,H,\eta)$. This is a generalisation of the description in the reductive case given in \cite{Nomizu1969} to the non-reductive case that appears, e.g., in~\cite{Figueroa-OFarrill2017_1}. We first define a (degenerate) symmetric bilinear form on $\fg$ by $\pair{X}{Y}=\eta\bigl(\overline{X},\overline{Y}\bigr)$ where $\overline{X}$, $\overline{Y}$ denote the images in $V=\fg/\fh$ of $X,Y\in\fg$; note that $\pair{-}{-}$ is $H$-invariant since $\eta$ is. We then define $U\colon \Odot^2\fg\to V$ by the equation
$2\eta(U(X,Y),v) = \pair{X}{\comm{Z}{Y}} + \pair{\comm{Z}{X}}{Y}
$
for all $X,Y\in\fg$, $v\in V$ and $Z\in\fg$ an arbitrary element such that $\overline{Z}=v$. This is well-defined by $\fh$-invariance of $\pair{-}{-}$. We can then define $\widetilde{L}\colon \fg\to\Hom(\fg,V)$ by
\[
	\widetilde{L}(X)Y = \frac{1}{2}\overline{\comm{X}{Y}} + U(X,Y).
\]
One can check that $\widetilde{L}(X)(Y)=0$ for all $X\in\fg$ and $Y\in\fh$, so $\widetilde{L}$ factors through a map $L\colon \fg\to\fgl(V)$ which is a Nomizu map and also satisfies
\begin{enumerate}\itemsep=0pt
	\item $\Im L\in\fso(V)$,
	\item $\widetilde{L}(X)(Y) - \widetilde{L}(X)(Y) - \overline{\comm{X}{Y}} = 0$ for all $X,Y\in \fg$.
\end{enumerate}
These two properties precisely say that the affine connection corresponding to $L$ is metric-compatible and torsion-free, hence it is the Levi-Civita connection.
